% !PRINTED 558
\weberpage{558}
%<<CONT

\gesperrt{II. Jede Zahl in $R(r)$, die die Substitution $(r, r_e)$ gestattet, also auch jede Zahl des Körpers $R(\eta)$ lässt sich als homogene lineare Function der Perioden $\eta, \eta_1 \ldots \eta_{e-1}$ darstellen.}

Denn nach I. können wir jede Zahl von $R(r)$ auf eine Weise in die Form

\weq{13}{\varphi(r) = \sum a_h r_h}

setzen, wo die Coëfficienten $a_h$ rationale Zahlen sind, und $h$ ein volles Restsystem nach dem Modul $n - 1$ durchläuft. Es ist aber
\[
\varphi(r_e) = \sum a_h r_{h+e} = \sum a_{h-e} r_h,
\]
und daher muss, wenn $\varphi(r) = \varphi(r_e)$ sein soll, $a_{h-e} = a_h$ sein für jedes $h$; also auch $a_h = a_{h+e} = a_{h+2e} \ldots$ Danach lässt sich (13) in die Form setzen
\[
\varphi(r) = \sum_{0,\,e-1}^{h} a_h r_h + \sum_{0,\,e-1}^{h} a_h r_{h+e} + \cdots + \sum_{0,\,e-1}^{h} a_h r_{h+(f-1)e},
\]
oder

\weq{14}{\varphi(r) = a\eta + a_1\eta_1 + \cdots + a_{e-1}\eta_{e-1},}

was zu beweisen war. Zu bemerken ist dabei noch, dass, wenn $\varphi$ ganzzahlige Coëfficienten hat, auch die Coëfficienten der $\eta_h$ ganze Zahlen werden.

Da jede Zahl des Körpers $R(\eta)$ auf eine und nur auf eine Art in der linearen Form (14) dargestellt werden kann, so nennen wir das Grössensystem $\eta, \eta_1, \eta_2 \ldots \eta_{e-1}$ eine \gesperrt{Basis des Körpers $R(\eta)$}.

Wenn man die Darstellung des Productes je zweier Perioden $\eta_h \eta_k$ in der Form (14) kennt, so kann man daraus durch eine wiederholte Anwendung alle rationalen Functionen der Perioden $\eta$ in derselben Form darstellen. Ebenso kann man auch leicht die Gleichung bilden, deren Wurzeln die $e$ Grössen $\eta$ sind, wenn man das Verfahren auf das Product

\weq{15}{F_e(x) = (x - \eta)(x - \eta_1) \cdots (x - \eta_{e-1})}

anwendet. Diese Gleichung lässt sich auch in Determinantenform darstellen.

Man setze nämlich

\weq{16}{\eta\eta_h = a_{0,h}\eta + a_{1,h}\eta_1 + \cdots + a_{e-1,h}\eta_{e-1},}

worin die $a_{0,h}, a_{1,h} \ldots a_{e-1,h}$ ganze Zahlen sind, über deren Berechnung später das Nähere folgen wird. Stellt man diese Gleichung für $h = 0, 1, \ldots e - 1$ in der Form auf:
%CONT>>
